\shortdefinition[Konfidenzintervall (K.I.)] mit Niveau $1 - \alpha$ ist Zufallsintervall $I = [A, B]$ s.d. $\forall \vartheta \in \Theta$ gilt:
\[
	\P_\vartheta[[A, B] \ni \vartheta] = \P_\vartheta[A \leq \vartheta \leq B] \geq 1 - \alpha
\]
mit $A, B$ Z.V. wie $A = a(\cX_1, \ldots, \cX_n)$, mit $a, b: \R^n \rightarrow \R$.
Dabei ist $\vartheta$ deterministisch. (\hl{Berechnung}: s. \ref{sec:distribution-function})

\shortremark[Approx. K.I.] Nur wenn genaue Verteilungsaussagen zur Verfügung stehen, sind exakte K.I. möglich,
sonst mittels zentralem Grenzwertsatz approx. berechnen ($z = z_{1 - \frac{\alpha}{2}}$)
\[
	P_\vartheta[|S_n^*| \leq z] \approx 1 - \alpha
\]
Auflösen von untenstehendem nach $\vartheta$ ($S_n$ wie in \ref{sec:weak-law-of-large-numbers})
\[
	|S_n - n \vartheta| \leq z \sqrt{n \vartheta(1 - \vartheta)} \text{ bzw. } (S_n - n \vartheta)^2 \leq z^2 n \vartheta(1 - \vartheta)
\]
Ist kompliziert, ersetzen der Ugl. durch Gl. (quad. Gl.)
\[
	\hat{\vartheta}_\pm = \frac{2 S_n + z^2 \pm \sqrt{(2 S_n + z^2)^2 - 4S_n^2(1 + \frac{z^2}{n})}}{2n(1 + \frac{z^2}{n})}
\]
Konfidenzintervall ist $[\hat{\vartheta}_-, \hat{\vartheta}_+]$.

\shade{gray}{Alternative Methoden} (Annahme $\vartheta(1 - \vartheta) \approx \frac{1}{4}$):
\[
	\left[\overline{S}_n - \frac{z}{2 \sqrt{n}}, \overline{S}_n + \frac{z}{2 \sqrt{n}}\right]
\]
(Ersetzen von $\vartheta$ durch $\overline{S}_n$):
\[
	\left[ \overline{S}_n - \frac{z}{2 \sqrt{n}} \sqrt{\overline{S}_n (1 - \overline{S}_n)}, \overline{S}_n + \frac{z}{2 \sqrt{n}} \sqrt{\overline{S}_n (1 - \overline{S}_n)} \right]
\]
Alle Methoden liefern unterschiedliche Resultate
